\documentclass[12pt,a4paper]{article}

% ------------------------------------------------
% Pacotes
% ------------------------------------------------
\usepackage[utf8]{inputenc}
\usepackage[T1]{fontenc}
\usepackage[brazil]{babel}
\usepackage{graphicx}
\usepackage[table]{xcolor}
\usepackage{array}
\usepackage{tabularx}
\usepackage{geometry}
\usepackage{fancyhdr}
\usepackage{mathptmx} % Times New Roman
\usepackage{enumitem}\usepackage{amsfonts}
\usepackage{amsmath} % Recomendado para o comando \text{}
\usepackage{amssymb}
\usepackage{tikz}
\usetikzlibrary{positioning, shapes.geometric}
\usepackage{listings}

\lstset{
    language=Java,
    basicstyle=\small\ttfamily,
    keywordstyle=\bfseries,
    showstringspaces=false,
    breaklines=true,
    frame=none,
    numbers=none,
    tabsize=2,
    commentstyle=\itshape\color{gray},
    morekeywords={this}
}

% TikZ Styles for ER Diagrams
\tikzset{
    block/.style={rectangle, draw, thick, minimum width=2cm, minimum height=1cm, fill=white},
    diamond_node/.style={diamond, draw, thick, minimum width=2cm, minimum height=1cm, fill=white, aspect=2},
    attribute/.style={ellipse, draw, thin, minimum width=1.5cm, minimum height=0.6cm, fill=white, font=\scriptsize},
    line/.style={draw, thick}, 
}
% ------------------------------------------------ 
% Configuração de página
% ------------------------------------------------
\geometry{
    a4paper,
    left=2cm,
    right=2cm,
    top=2cm,
    bottom=2cm,
    headheight=15pt % Adicione isso
}% Caixa para resposta discursiva
\newcommand{\resposta}[1]{
\noindent
\framebox[\textwidth]{
\begin{minipage}{0.98\textwidth}
\vspace{#1}
\end{minipage}
}
}

\pagestyle{fancy}
\fancyhf{}
\renewcommand{\headrulewidth}{0pt}

\setlength{\parindent}{0pt}

% ------------------------------------------------
% Comandos auxiliares
% ------------------------------------------------
\newcolumntype{L}[1]{>{\raggedright\arraybackslash}p{#1}}
\newcolumntype{C}[1]{>{\centering\arraybackslash}p{#1}}

% ------------------------------------------------
% Documento
% ------------------------------------------------
\begin{document}

% ------------------------------------------------
% Cabeçalho - Corrigido para evitar aninhamento problemático
% ------------------------------------------------
\begin{tabularx}{\textwidth}{@{} L{4cm} X @{}}
    \raisebox{-0.5cm}{\includegraphics[width=4cm]{Universidade_de_Rio_Verde.jpg}} &
    \hfill
    \begin{tabular}{|>{\centering\arraybackslash}m{10cm}|>{\raggedright\arraybackslash}p{1.8cm}|}
        \hline
        \cellcolor{black}\color{white}\rule{0pt}{1.2cm}\textbf{AVALIAÇÃO 1 (ADAPTADA)} 
        & 
        \scriptsize Nota \tabularnewline
        \hline
    \end{tabular} 
\end{tabularx}
\vspace{0.3cm}

% ------------------------------------------------  
% Tabela principal - Estrutura limpa
% ------------------------------------------------
\begin{tabularx}{\textwidth}{|X|L{2cm}|}
\hline
\textbf{Curso} & \textbf{Campus} \\
\textit{Engenharia de Software} & \textit{Rio Verde} \\
\hline
\multicolumn{2}{|l|}{ \textit{Programação Orientada a Objetos}} \\
\hline
\textbf{Nome do(a) acadêmico(a)} & \textbf{Valor} \\
& 7,0 \\
\hline
\textbf{Professor} & \textbf{Turma} \\
\textit{ Sergio Souza Novak} & \textit{B} \\
\hline
\textbf{Data da Avaliação} & \textbf{Código} \\
\textit{20/03/2026} & \textit{ESW448} \\
\hline
\end{tabularx}
 
% ------------------------------------------------
% Barra de Atenção
% ------------------------------------------------
\vspace{0.2cm}

\begin{tabularx}{\textwidth}{|X|}
\hline
\cellcolor{black}\color{white}
\centering \textbf{ATENÇÃO:} Somente serão passíveis de \textbf{REVISÃO} avaliações resolvidas à \textbf{TINTA}. \tabularnewline
\hline
\end{tabularx}

\vspace{0.3cm}

\textbf{ORIENTAÇÕES PARA A RESOLUÇÃO}  

\begin{itemize}
    \item A avaliação é individual, sem consulta;
    \item A interpretação de texto faz parte da avaliação;
\end{itemize}

\vspace{0.5cm}

\begin{center}
    \rule{\textwidth}{0.4pt}
    \textbf{QUESTÕES}
\end{center}

\vspace{0.5cm} 

\begin{enumerate}
    \item \textbf{(ENADE 2017)}: O encapsulamento é um mecanismo da programação orientada a objetos no qual os membros de uma classe (atributos e métodos) constituem uma caixa preta. O nível de visibilidade dos membros pode ser definido pelos modificadores de visibilidade "privado", "público" e "protegido". Com relação ao comportamento gerado pelos modificadores de visibilidade, assinale a opção correta. 
    \begin{enumerate}[label=\Alph*)]
        \item Um atributo privado pode ser acessado pelos métodos privados da própria classe e pelos métodos protegidos das suas classes descendentes.
        \item Um atributo privado pode ser acessado pelos métodos públicos da própria classe e pelos métodos públicos das suas classes descendentes.
        \item Um membro público é visível na classe à qual ele pertence, mas não é visível nas suas classes descendentes.
        \item Um método protegido não pode acessar os atributos privados e declarados na própria classe.
        \item Um membro protegido é visível na classe à qual pertence e em suas classes descendentes.
    \end{enumerate}

    \item \textbf{Encapsulamento:} Observe o diagrama da classe \textbf{ContaBancaria} abaixo. Assinale a alternativa que apresenta a implementação correta em Java, respeitando o encapsulamento e a lógica dos métodos.

    \begin{center} 
    \begin{tikzpicture}[node distance=0cm, outer sep=0pt, every node/.style={draw, rectangle, minimum height=0.6cm, font=\small\ttfamily, anchor=north west}]
        \node (header) [fill=white, minimum width=9cm, font=\small\bfseries] at (0,0) {ContaBancaria};
        
        \node (attrs) [minimum width=9cm, align=left] at (0,-0.6cm) {
            - saldo : double \\
            + depositar(valor: double) : void \\
            + sacar(valor: double) : void \\
            + getSaldo() : double
        };
    \end{tikzpicture}
    \end{center}

    \begin{enumerate}[label=\Alph*)]
        \item \begin{lstlisting}
public class ContaBancaria {
    public double saldo;
    public void depositar(double v) { saldo = v; }
}
        \end{lstlisting}
        \item \begin{lstlisting}
public class ContaBancaria {
    private double saldo;
    public ContaBancaria(double s) { this.saldo = s; }
    public void depositar(double v) { this.saldo += v; }
    public void sacar(double v) { this.saldo -= v; }
    public double getSaldo() { return this.saldo; }
}
        \end{lstlisting}
        \item \begin{lstlisting}
public class ContaBancaria {
    private double saldo;
    public void depositar(double v) { }
    public void sacar(double v) { }
}
        \end{lstlisting}
        \item \begin{lstlisting}
public class ContaBancaria {
    private double saldo;
    private double getSaldo() { return saldo; }
}
        \end{lstlisting}
        \item \begin{lstlisting}
public class ContaBancaria {
    public double getSaldo() { return 100.0; }
}
        \end{lstlisting}
    \end{enumerate}

    \item Na programação orientada a objetos, os modificadores de acesso controlam a visibilidade de classes, atributos e métodos. Assinale a alternativa que apresenta o modificador que permite que um membro seja acessado apenas dentro da própria classe e por suas subclasses:
    \begin{enumerate}[label=\Alph*)]
        \item \texttt{public}
        \item \texttt{private}
        \item \texttt{protected}
        \item \texttt{package-private}
        \item \texttt{static}
    \end{enumerate}

    \item \textbf{Construtores:} Qual é a principal finalidade de um método construtor em uma classe?
    \begin{enumerate}[label=\Alph*)]
        \item Destruir o objeto quando ele não for mais necessário.
        \item Inicializar os atributos de um objeto no momento de sua criação.
        \item Permitir que o objeto seja herdado por outras classes.
        \item Converter um objeto em uma string para exibição.
        \item Impedir que o objeto seja clonado.
    \end{enumerate}


    \vspace{0.3cm}


    \item \textbf{Indústria 4.0 — Máquina CNC:} No contexto de sistemas de manufatura inteligentes, a classe \texttt{MaquinaCNC} foi projetada para representar uma máquina de usinagem com controle automatizado.

    \begin{center}
    \begin{tikzpicture}[node distance=0cm, outer sep=0pt, every node/.style={draw, rectangle, minimum height=0.6cm, font=\small\ttfamily, anchor=north west}]
        \node (header) [fill=white, minimum width=9cm, font=\small\bfseries] at (0,0) {MaquinaCNC};
        
        \node (attrs) [minimum width=9cm, align=left] at (0,-0.6cm) {
            - modelo : String \\
            - horasOperacao : double \\
            - temperaturaAtual : double \\
            - temperaturaMaxima : double \\
            - ligada : boolean \\
            - pecasProduzidas : int \\
            + MaquinaCNC(modelo: String, tempMax: double) \\
            + ligar() : void \\
            + desligar() : void \\
            + produzir(qtd:int) : boolean \\
            + getTemperaturaAtual() : double \\
            + precisaManutencao() : boolean \\
            - controlarTemperatura(qtd:int) : void \\
            - verificarManutencao() : boolean
        };
    \end{tikzpicture}
    \end{center}
\begin{small}
\begin{lstlisting}
01: // MaquinaCNC.java
02: public class MaquinaCNC {
03:     private String modelo;
04:     private double horasOperacao;
05:     private double temperaturaAtual;
06:     private double temperaturaMaxima;
07:     private boolean ligada;
08:     private int pecasProduzidas;
09: 
10:     public MaquinaCNC(String modelo, double tempMax) {
11:         this.modelo = modelo;
12:         this.temperaturaMaxima = tempMax;
13:         this.temperaturaAtual = 25.0; // temperatura ambiente
14:     }
15: 
16:     public void ligar() { this.ligada = true; }
17:     public void desligar() { this.ligada = false; }
18: 
19:     public boolean produzir(int qtd) {
20:         if (!this.ligada){ return false; }
21: 
22:         this.pecasProduzidas += qtd;
23:         this.horasOperacao += qtd * 0.1;
24:         this.controlarTemperatura(qtd);
25:         return true;
26:     }
27: 
28:     public double getTemperaturaAtual() { return this.temperaturaAtual; }
29:     public boolean precisaManutencao() { return this.verificarManutencao(); }
30: 
31:     private void controlarTemperatura(int qtd) {
32:         this.temperaturaAtual += qtd * 2.0;
33:         if (this.temperaturaAtual > temperaturaMaxima) {
34:             desligar();
35:         }
36:     }
37: 
38:     private boolean verificarManutencao() {
39:         return horasOperacao > 1000;
40:     }
41: }
42: 
43: // Main.java
44: public class Main {
45:     public static void main(String[] args) {
46:         MaquinaCNC cnc = new MaquinaCNC("CNC-X1", 100.0);
47:         cnc.ligar();
48:         if (cnc.produzir(50)) {
49:             System.out.println("Producao concluida.");
50:         }
51:     }
52: }
\end{lstlisting}
\end{small}

    Com base no código da classe \texttt{MaquinaCNC}, analise as afirmações abaixo como Verdadeiras (V) ou Falsas (F) e assinale a alternativa que apresenta a sequência \textbf{correta}:

    ( ) Se a linha 14 for removida, a máquina poderá produzir peças (\texttt{pecasProduzidas}) mesmo estando desligada (\texttt{ligada = false}). \\
    ( ) Os atributos das linhas 03 a 08 estão corretamente encapsulados como \texttt{private}, impedindo o acesso direto de outras classes. \\
    ( ) O método \texttt{precisaManutencao} (linha 22) expõe diretamente o atributo \texttt{horasOperacao}, violando o princípio do encapsulamento.

    \begin{enumerate}[label=\Alph*)]
        \item V - V - V - F.
        \item V - F - V - V.
        \item F - V - F - V.
        \item V - V - F - F.
        \item F - F - V - V.
    \end{enumerate}



    \item \textbf{Siderurgia — Alto Forno:} A modelagem da classe \texttt{AltoForno} envolve o controle de um equipamento utilizado para produção de metal.

    \begin{center}
    \begin{tikzpicture}[node distance=0cm, outer sep=0pt, every node/.style={draw, rectangle, minimum height=0.6cm, font=\small\ttfamily, anchor=north west}]
        \node (header) [fill=white, minimum width=9cm, font=\small\bfseries] at (0,0) {AltoForno};
        
        \node (attrs) [minimum width=9cm, align=left] at (0,-0.6cm) {
            - temperaturaInterna : double \\
            - temperaturaMinimaOperacao : double \\
            - temperaturaMaximaOperacao : double \\
            - volumeMetalProduzido : double \\
            - ativo : boolean \\
            + AltoForno(tempMin: double, tempMax: double) \\
            + iniciarForno() : void \\
            + desligarForno() : void \\
            + injetarCombustivel(qtd:double) : double \\
            + produzirMetal() : double \\
            + estaAtivo() : boolean \\
            - verificarFaixaTermica() : boolean \\
            - controleSeguranca() : void
        };
    \end{tikzpicture}
    \end{center}
\begin{small}
\begin{lstlisting}
01: // AltoForno.java
02: public class AltoForno {
03:     private double temperaturaInterna;
04:     private double temperaturaMinimaOperacao;
05:     private double temperaturaMaximaOperacao;
06:     private double volumeMetalProduzido;
07:     private boolean ativo;
08: 
09:     public AltoForno(double tempMin, double tempMax) {
10:         this.temperaturaMinimaOperacao = tempMin;
11:         this.temperaturaMaximaOperacao = tempMax;
12:         this.temperaturaInterna = 20.0;
13:     }
14: 
15:     public void iniciarForno() { this.ativo = true; }
16:     public void desligarForno() { this.ativo = false; }
17:     public boolean estaAtivo() { return ativo; }
18: 
19:     public double injetarCombustivel(double qtd) {
20:         this.temperaturaInterna += qtd * 10;
21:         controleSeguranca();
22:         return temperaturaInterna;
23:     }
24: 
25:     public double produzirMetal() {
26:         if (ativo && verificarFaixaTermica()) {
27:             double producao = 100.0;
28:             this.volumeMetalProduzido += producao;
29:             return producao;
30:         }
31:         return 0.0;
32:     }
33: 
34:     private boolean verificarFaixaTermica() {
35:         return temperaturaInterna >= temperaturaMinimaOperacao && 
36:                temperaturaInterna <= temperaturaMaximaOperacao;
37:     }
38: 
39:     private void controleSeguranca() {
40:         if (temperaturaInterna > temperaturaMaximaOperacao) {
41:             desligarForno();
42:         }
43:     }
44: }
\end{lstlisting}
\end{small}

    Considerando o código da classe \texttt{AltoForno} e os princípios de Programação Orientada a Objetos, julgue os itens a seguir como Verdadeiras (V) ou Falsas (F) e assinale a sequência \textbf{correta}:

    ( ) O método \texttt{verificarFaixaTermica} (linhas 28-31) é \texttt{private}, o que significa que ele só pode ser invocado por outros métodos dentro de outras classes. \\
    ( ) Na linha 15, o método \texttt{controleSeguranca} é chamado para garantir que a temperatura não ultrapasse o limite após a injeção de combustível, demonstrando o controle de estado interno. \\
    ( ) A lógica de produção de metal na linha 20 depende exclusivamente do forno estar \texttt{ativo}, ignorando a temperatura atual do equipamento. \\
    ( ) O atributo \texttt{temperaturaInterna} (linha 03) poderia ser \texttt{public} sem afetar a segurança do equipamento, pois o método \texttt{controleSeguranca} corrigiria qualquer valor inválido.

    \begin{enumerate}[label=\Alph*)]
        \item V - V - F - F.
        \item V - F - V - V.
        \item F - V - V - F.
        \item V - V - V - F.
        \item F - F - F - V.
    \end{enumerate}


    \item \textbf{(VUNESP 2021 - Semae de Piracicaba - SP - Programador Junior)} Em programação orientada a objetos, métodos de acesso do tipo setter têm a finalidade primária de:
    \begin{enumerate}[label=\Alph*)]
        \item obter o valor de um atributo.
        \item modificar o valor de um atributo.
        \item inicializar os valores de atributos de um objeto assim que ele é instanciado.
        \item contar quantas vezes o valor de um atributo foi obtido.
        \item salvar o valor de um atributo de um objeto em um banco de dados.
    \end{enumerate}

    \item \textbf{(FUNDATEC 2025 - GHC-RS - Programador)} Considerando os conceitos de orientação a objetos na linguagem de programação Java, assinale a alternativa que apresenta quem pode acessar o atributo \texttt{v} diretamente:

\begin{lstlisting}
package conceicao; 
public class Conceicao {       
    protected String v; 
}
\end{lstlisting}

    \begin{enumerate}[label=\Alph*)]
        \item Apenas as classes que estejam no mesmo pacote.
        \item Apenas os métodos da própria classe.
        \item As classes filhas de qualquer pacote.
        \item As classes filhas apenas do mesmo pacote.
        \item Por qualquer classe de qualquer pacote.
    \end{enumerate}

    \item Considere a modelagem de uma classe \texttt{Carro} que possui um atributo privado chamado \texttt{velocidadeAtual}. Sobre a importância de se utilizar métodos \textit{getters} e \textit{setters}, assinale a alternativa que descreve corretamente uma vantagem e um exemplo de regra de negócio que poderia ser implementada:

    \begin{center}
    \begin{tikzpicture}[node distance=0cm, outer sep=0pt, every node/.style={draw, rectangle, minimum height=0.6cm, font=\small\ttfamily, anchor=north west}]
        \node (header) [fill=white, minimum width=9cm, font=\small\bfseries] at (0,0) {Carro};
        
        \node (attrs) [minimum width=9cm, align=left] at (0,-0.6cm) {
            - velocidadeAtual : double \\
            + getVelocidade() : double \\
            + setVelocidade(v: double) : void
        };
    \end{tikzpicture}
    \end{center}

    \begin{enumerate}[label=\Alph*)]
        \item A vantagem é permitir que qualquer pessoa altere a velocidade sem restrições. Exemplo: Deixar a velocidade ser negativa.
        \item A vantagem é o encapsulamento, permitindo validar os dados antes de alterar o atributo. Exemplo: Impedir que a velocidade receba valores negativos.
        \item A vantagem é diminuir o tamanho do código Java. Exemplo: Somar a velocidade com o ano do carro.
        \item A vantagem é tornar todos os atributos públicos automaticamente. Exemplo: Alterar a cor do carro no método da velocidade.
        \item A vantagem é que o código roda mais rápido. Exemplo: Multiplicar a velocidade por zero para parar o carro.
    \end{enumerate}

\end{enumerate}
\vspace{0.8cm}


\end{document}
